\documentclass[11pt]{article}
\usepackage[margin=1in]{geometry}
\usepackage{amsmath,amssymb,amsthm}
\usepackage[T1]{fontenc}
\usepackage[utf8]{inputenc}
\usepackage{lmodern}
\usepackage{graphicx}
\usepackage[hidelinks]{hyperref}
\emergencystretch=2em

\newtheorem{theorem}{Theorem}
\newtheorem{lemma}{Lemma}
\newtheorem{corollary}{Corollary}

\title{Bounded-Excess Cotype for Gaussian Random Polytopes}
\author{Run \texttt{fa\_banach\_001}, agent \texttt{agent\_lane\_15}}
\date{2026-06-29}

\begin{document}
\maketitle

\section*{Status}

\textbf{Status:} partial result, likely valid.  This solves the bounded
vertex-excess subcase $N-n\le M$ of the arbitrary-$N$ cotype problem in
arXiv:2603.04749, with constants depending on $M$.

\section*{Source Problem}

Huang and Tikhomirov prove in \cite{HT} that, when
$K'\ge N/n\ge K>1$, the Gaussian polytope
\[
  P_{N,n}=\operatorname{conv}\{\pm X_i:1\le i\le N\}\subset \mathbb R^n
\]
generates, with high probability, a normed space
$(\mathbb R^n,\|\cdot\|_{P_{N,n}})$ of finite dimension-independent cotype.
Their Open Problems section asks whether there are universal
$q,C_q<\infty$ such that the same conclusion holds for every $N\ge n$.
The same section asks the analogous question for random orthogonal
projections of the standard cross-polytope.

\begin{center}
\includegraphics[width=\linewidth]{figures/open_problem_crop.png}
\end{center}

\section*{Partial Result}

\begin{theorem}[bounded excess]
Let $M\ge 0$ be fixed.  Let $N\ge n$, let $X_1,\ldots,X_N$ span
$\mathbb R^n$, and assume $N-n\le M$.  Then
$(\mathbb R^n,\|\cdot\|_{P_{N,n}})$ has cotype $2$ with a constant depending
only on $M$.  In particular, for Gaussian $X_i$ this holds almost surely.
\end{theorem}

\begin{corollary}[bounded-codimension projected cross-polytopes]
Let $M\ge 0$ be fixed.  If $L\subset\mathbb R^N$ is any $n$-dimensional
subspace with $N-n\le M$, then the orthogonal projection of the standard
cross-polytope $B_1^N$ onto $L$ defines on $L$ a normed space of cotype $2$
with constant depending only on $M$.
\end{corollary}

\section*{Idea of the Proof}

The missing near-square case is not random at bounded excess.  The Gaussian
polytope is the image of the $\ell_1^N$ unit ball under the map
$A e_i=X_i$.  Thus the associated normed space is the quotient
$\ell_1^N/\ker A$.  When $N-n$ is bounded, this quotient is by a
bounded-dimensional subspace.  Auerbach's lemma gives a projection onto that
kernel whose norm is controlled by its dimension, so the quotient is uniformly
isomorphic to a subspace of $\ell_1^N$.  The cotype $2$ of $\ell_1$ then
transfers to the quotient.

\section*{Proof}

We use two standard facts.  First, $\ell_1$ has cotype $2$; denote one
admissible cotype constant by $C_{\ell_1}$.  Second, if $Y$ is $D$-isomorphic
to a subspace of a space of cotype $q$ with constant $C_q$, then $Y$ has
cotype $q$ with constant at most $D C_q$.

\begin{lemma}[finite-dimensional quotients]
Let $X$ be a Banach space of cotype $q$ with constant $C_q$, and let
$F\subset X$ be an $m$-dimensional subspace.  Then $X/F$ has cotype $q$ with
constant at most $(m+1)C_q$.
\end{lemma}

\begin{proof}
If $m=0$ there is nothing to prove.  For $m\ge 1$, choose an Auerbach basis
$(u_j,u_j^*)_{j=1}^m$ of $F$, so $\|u_j\|=\|u_j^*\|=1$ and
$u_i^*(u_j)=\delta_{ij}$.  Extend each $u_j^*$ to a norm-one functional
$\widetilde u_j^*\in X^*$ by Hahn--Banach, and define
\[
  P x=\sum_{j=1}^m \widetilde u_j^*(x)u_j .
\]
Then $P$ is a projection from $X$ onto $F$ and $\|P\|\le m$.  The map
\[
  T:X/F\to \ker P,\qquad T(x+F)=(I-P)x,
\]
is well-defined, linear, and bijective.  Moreover
$\|T\|\le \|I-P\|\le m+1$, while the inverse map $\ker P\to X/F$ has norm at
most $1$, because the quotient norm never exceeds the norm of a representative.
Thus $X/F$ is $(m+1)$-isomorphic to the subspace $\ker P$ of $X$.  Cotype
passes to subspaces and then across this isomorphism with the stated
distortion loss.
\end{proof}

Now set $A:\ell_1^N\to\mathbb R^n$ by $Ae_i=X_i$.  Since the $X_i$ span
$\mathbb R^n$, the map $A$ is onto and
\[
  A(B_1^N)=\operatorname{conv}\{\pm X_i:1\le i\le N\}=P_{N,n}.
\]
Therefore, for $x\in\mathbb R^n$,
\[
  \|x\|_{P_{N,n}}
  =\inf\{\|a\|_1:a\in\mathbb R^N,\ Aa=x\},
\]
which is exactly the quotient norm on $\ell_1^N/\ker A$.  Since
\[
  \dim \ker A=N-n\le M,
\]
the lemma applied to $X=\ell_1^N$ gives cotype $2$ with constant at most
$(M+1)C_{\ell_1}$.  For independent standard Gaussian vectors, spanning
$\mathbb R^n$ holds almost surely when $N\ge n$.

The projected cross-polytope statement is identical: an orthogonal projection
$Q:\ell_1^N\to L$ maps $B_1^N$ onto the projected cross-polytope, and the
associated normed space is $\ell_1^N/\ker Q$ with
$\dim\ker Q=N-n\le M$.

\section*{Limitations}

The constants depend on the excess $M=N-n$.  Hence this does not settle the
full arbitrary-$N$ problem, which asks for constants uniform over all
$N\ge n$.  It also does not cover regimes where $N-n\to\infty$, including the
large-aspect random-projection problem isolated in the source paper and in the
previous attempt note for this target.

\section*{Novelty and Duplicate Check}

The run indexes contained one previous attempt for arXiv:2603.04749 and no
solution packet for this bounded-excess subcase.  Web searches on
2026-06-29 for the exact arXiv id and title, for bounded-codimension/random
polytope cotype phrases, and for finite-dimensional quotient/cotype variants
found the arXiv source page and standard Auerbach projection facts, but no
direct later solution of this subcase.  The proof above is deterministic and
uses only standard Banach-space facts.

\section*{Human Review Recommendation}

Review the finite-dimensional quotient lemma and the identification of
$P_{N,n}$ as a quotient image of $B_1^N$.  If those are accepted, the bounded
excess partial result is complete.

\begin{thebibliography}{9}
\bibitem{HT}
H. Huang and K. Tikhomirov,
\emph{Cotype of random polytopes},
arXiv:2603.04749, 2026.

\bibitem{MP}
B. Maurey and G. Pisier,
\emph{Series de variables aleatoires vectorielles independantes et proprietes
geometriques des espaces de Banach},
Studia Math. 58 (1976), 45--90.

\bibitem{LT}
J. Lindenstrauss and L. Tzafriri,
\emph{Classical Banach Spaces I and II},
Springer, 1996.
\end{thebibliography}

\end{document}
